\section{Prefetch-Analyse}

Windows Prefetch: Dateien im Verzeichnis \texttt{Windows/Prefetch/} mit Endung \texttt{.pf}.
Enthalten Informationen über ausgeführte Programme (Name, Run Count, letzte Ausführungszeit).
Prefetch lädt häufig verwendete DLLs vor — zeigt auch Autostart-Malware.

\begin{lstlisting}[language=bash]
# Prefetch-Dateien finden
fls -pr /dev/loopXpN | grep -i prefetch
mount -o ro /dev/loopXpN /mnt
find /mnt -type f -iname "*.pf" -print0 | xargs -0 ls -la

# Prefetch analysieren
sccainfo /mnt/Windows/Prefetch/PROGRAMM.EXE-HASH.pf
# Ausgabe: Run Count, Last Run Time(s)
\end{lstlisting}

\section{Autostart-Einträge}

Registry Run-Keys zeigen Programme, die bei jedem Login starten:

\begin{lstlisting}[language=bash]
regripper -p run -r /tmp/SOFTWARE
regripper -p run -r /tmp/NTUSER.DAT
\end{lstlisting}

Geplante Tasks:

\begin{lstlisting}[language=bash]
fls -pr /dev/loopXpN | grep -i 'Tasks'
icat /dev/loopXpN <INODE_TASK> > /tmp/task.xml
cat /tmp/task.xml
\end{lstlisting}

\section{YARA-Regeln}

Pattern-Matching für Malware-Erkennung.
Regelformat:

\begin{lstlisting}
rule rule_name {
    strings:
        $var1 = "string" nocase
        $var2 = "string2" nocase
    condition:
        $var1 and $var2
}
\end{lstlisting}

Beispiel — UPX-gepackte Binaries erkennen:

\begin{lstlisting}
rule upx_find {
    strings:
        $upx0 = "upx0" nocase
        $upx1 = "upx1" nocase
        $rsrc = ".rsrc" nocase
    condition:
        $upx0 and $upx1 and $rsrc
}
\end{lstlisting}

Beispiel — Gestohlene Signaturen (Stuxnet-Referenz):

\begin{lstlisting}
rule stolen_signatures {
    strings:
        $find1 = "realtek" nocase
        $find2 = "jmicron" nocase
    condition:
        $find1 or $find2
}
\end{lstlisting}

Anwendung: \texttt{yara rules.yar suspicious\_file}

\section{Event Logs (EVTX)}

Windows Event Logs für Security-Analyse:

\begin{lstlisting}[language=bash]
fls -pr /dev/loopXpN | grep -i '\.evtx$'
icat /dev/loopXpN <INODE> > /tmp/security.evtx
evtx_dump.py /tmp/security.evtx | grep -i "malware\|virus\|threat"
\end{lstlisting}

\section{Weitere Indikatoren}

\begin{description}
	\item[Hosts-Datei:] DNS-Hijacking erkennen (\texttt{icat} → \texttt{hosts} prüfen)
	\item[Strings-Analyse:] Verdächtige Binaries nach Schlüsselwörtern durchsuchen:
\end{description}

\begin{lstlisting}[language=bash]
icat /dev/loopXpN <INODE_EXE> > /tmp/suspicious.exe
strings /tmp/suspicious.exe | grep -Ei "http|cmd|powershell|download|bitcoin"
\end{lstlisting}
